\section{引言}

设想小明正在分阶段减烟方案之间进行选择。两个方案可以具有相同的初始状态、终止状态、持续时间和累计暴露指标，却在是否反向、跨越哪些状态水平以及是否到达治疗目标等方面不同。本文的问题是：对于一个可加的通行评价而言，这类历史中的哪些信息可以无损压缩？这种压缩又会对可观察选择施加什么限制？

一条历史首先生成一个\emph{通行账户}。该账户记录所处一侧、方向、跨越的状态水平、重复次数以及目标达成。它有意忽略日历时间，也忽略那些在不改变有向跨越记录的情况下可以交换的闭合游程之先后顺序。这种抽象不是直接假定的，而是被刻画出来的：任何对弱时间变换不变且对路径拼接可加的赋值，都必须通过一个通用载体分解；定理~\ref{thm:fiber} 则证明其逆向的纤维刻画。两条有限步历史具有相同账户，当且仅当它们可以通过可行的时间变换、停留、同基点闭合游程交换以及闭合路径的循环旋转相互连接。因此，该定理精确说明了账户究竟遗忘了哪些时间顺序信息。

一维几何还带来第二个压缩结果。净有向通行由有序端点决定，因此有序端点与总跨越剖面共同即可重构完整的有向账本，包括目标处的轨迹信息。这给出一个更小的充分记录，并说明：一旦同时观察到端点流量与总通行，方向就不再是独立的原始量。

通行账户本身不包含偏好。可加测量从实际历史之间的比较出发，而这个实际历史域并不对任意形式加法封闭。平衡通行循环一致性（BPCC）排除这样一类平衡的揭示改进列表：其中账户差分总和恰好抵消但至少有一项严格改进；由此得到一个精确的偏可加次序。成对消失基准一致性（PVBC）进一步排除相对于一个固定全支撑基准呈无穷小的实际严格比较。它使实际历史次序可以由一族归一化实数支撑赋值的一致同意来精确表示，同时仍允许未揭示的形式方向保持偏序。

正则性是对上述表示的进一步细化。穿孔区间 $(a,b_n]$ 在 $b_n\downarrow a$ 时从通行记录中消失，而区间 $[0,b_n]$ 始终包含一次目标事件。穿孔揭示次序闭性（PROC）要求揭示比较对第一类消失扰动保持闭合。于是，正则支撑赋值诱导出四个有向 Stieltjes 测度，并允许在目标处存在一个原子。面向目标的方向单调性则为四个模块提供应用层面的方向性解释。

对于每个正则支撑赋值 $\Lambda$，一维跨越守恒给出
\[
 U_\Lambda(\gamma)
 =P_\Lambda(\gamma(1))-P_\Lambda(\gamma(0))
  +\sum_s\int V_s^\gamma\,\dd\kappa_{s,\Lambda}
  +\sum_s\alpha_{s,\Lambda}A_s(\gamma).
\]
这一分解主要用于解释：端点进展被从对方向反转保持偶对称的总通行部分以及目标达成部分中分离出来。当多个支撑赋值都与原始次序相容时，这些分量依赖于具体赋值，因此可观察结论应以序数形式或集合识别形式表述。

通行账户的经济学用途并不要求所有背景后果在任意历史之间都字面不变。命题~\ref{prop:matched-design} 构造了一类匹配方案的时间安排，而无需允许任意快的遍历。对于任意有限个具有共同端点、且存在应用上可行计时遍历的右侧方案，只要共同时间跨度足够长，允许在端点停留就足以同时匹配持续时间以及一个在两端点取不同值的连续流量统计量的累计值，同时保持每条实际遍历及其通行账户不变。因此，可以通过实验设计固定端点、持续时间以及某个累计暴露量，而让通行几何发生变化。

由此得到两个可观察含义。第一，相同的通行账户差分可以把原始序数排名跨方案环境转移。第二，在有限经验网格上，有限个已观察选择定义了一个可加补全多面体，其线性像为任意未观察通行对比给出锐界。货币计价物可以把同一有限设计中的界校准为货币单位，但真正的识别对象是与数据相容的通行对比区间，而不是人为选定的一个标量效用。

数学工具主要来自可加测量、有序向量空间分离、定性概率与 Stieltjes 表示。本文与具体路径域相关的结果包括：精确路径纤维；由端点和总通行在一维中重构方向；在非加法封闭的历史域上进行可加定价；消失通行与保留目标达成之间的目标边界；以及即使只涉及六条历史，其系数重复次数也可能无界、但仍可由实际路径实现的消去证书。减烟是贯穿全文的方案环境；形式结果适用于具有相同通行语义的一维目标相对调整问题。

全文安排如下。第~\ref{sec:recovery}--\ref{sec:environment} 节定义经济对象与路径域。第~\ref{sec:universal} 节构造通用账户、精确纤维与一维压缩。第~\ref{sec:ordinal}--\ref{sec:stieltjes} 节建立可加支撑赋值以及正则测度化细化。第~\ref{sec:pha} 节给出端点--总通行--目标达成的解释和方案转移限制。第~\ref{sec:implications} 节讨论匹配设计与有限数据预测。附录集中给出纤维定理证明、分离论证、正则性构造及消去复杂度结果。
